
cet
élément du @@M0@@ est-il indépendant du choix particulier des @@M1@@ d'intersection @@M2@@?

\section*{7. Classes de Chern des complexes parfaits}

\textbf{7.1.} Lorsqu'on se pose la question de développer une formule de Riemann--Roch cohomologique, avec ou sans dénominateurs, en dehors de l'hypothèse assez artificielle d'existence d'un Module inversible ample sur les schémas envisagés, on est nécessairement confronté à la question d'une définition, pour un complexe de Modules parfait @@M3@@ sur un schéma @@M4@@, des classes de Chern
@@M5@@
généralisant les classes de Chern des Modules localement libres SGA 5 VII. La même question se pose d'ailleurs pour des topos localement annelés généraux (comparer [10]), ainsi que pour les complexes parfaits sur des espaces paracompacts annelés par les fonctions complexes continues, en travaillant alors avec la cohomologie entière ordinaire. Dans le cas d'un espace compact, on sait que le @@M6@@ construit avec les complexes parfaits est le même que le ``naif'' construit avec les Modules localement libres (Exp. IV), et dans ce cas on est ramené trivialement à la définition classique. Mais déjà dans le cas d'un espace localement compact le problème posé semble non trivial, puisque la cohomologie entière d'un tel espace n'est pas toujours la limite projective des cohomologies de ses parties compactes.

\textbf{7.2.} Plus précisément, lorsque @@M7@@ est une partie fermée de @@M8@@ telle que @@M9@@ soit nul (au sens des catégories dérivées), i.e. telle que les faisceaux de cohomologie de @@M10@@ soient à supports dans @@M11@@, on aimerait définir des classes de Chern \uline{à support dans} @@M12@@
@@M13@@
qui redonneraient les classes de Chern ``ordinaires'' (sic) de 7.1 à l'aide de l'homomorphisme canonique
@@M14@@
